\subsection{Requirements}
\label{subsec:requirements}
Deelvraag 1 vraagt aan welke eisen een risicoscore moet voldoen om gebruikers te ondersteunen bij het beoordelen van e-mails. Het effect op gebruikers valt buiten de scope van dit onderzoek. De eisen worden afgeleid uit sectie~\ref{subsubsec:uitleg-voor-gebruikers}.

\begin{description}
  \item[R1: weinig \emph{false positives}.] Gebruikers wennen aan foutieve waarschuwingen die ze herhaaldelijk zien \autocite{Anderson2015}. False positives moeten daarom zo veel mogelijk vermeden worden.
  \item[R2: aanwijzingen bij de bron.] Waarschuwingen werken beter dicht bij de verdachte link \autocite{Petelka2019}. De uitleg moet gebruikers dus kunnen wijzen op de plaats in de e-mail waar de aanwijzing staat.
  \item[R3: begrijpelijke taal.] Veel gebruikers begrijpen technische waarschuwingen niet \autocite{Egelman2008}. De uitleg moet dus in gewone taal en zonder jargon zijn geschreven.
  \item[R4: geen overdreven zekerheid.] Een overtuigende uitleg kan tot \emph{overreliance} leiden \autocite{Bucinca2021}. De uitleg mag dus geen overdreven zekerheid uitstralen en moet overeenkomen met de werkelijke beoordeling.
\end{description}